\documentclass[10pt,a4paper]{amsart}
\usepackage[margin=19mm]{geometry}
\usepackage{amsmath,amssymb,mathtools}
\usepackage[scr=rsfs,cal=euler]{mathalfa}
\usepackage{xfrac}
\usepackage[unicode,hidelinks,bookmarks=false]{hyperref}
\newcommand{\ie}{i.e.\ }
\newcommand{\cf}{cf.\ }
\newcommand{\resp}{respectively\ }
\newcommand{\Subsection}{\subsection}
\providecommand{\nobreakdash}{\nobreak\hbox{-}}
\setlength{\emergencystretch}{2em}
\multlinegap0pt
\usepackage{bm}
\usepackage{bbm}
\usepackage{dsfont}
\usepackage{xspace}
\usepackage{cancel}
\usepackage{mathtools}
\usepackage{amsthm}
\makeatletter
\@ifundefined{claim}{\newtheorem{claim}{Claim}}{}
\@ifundefined{theorem}{\newtheorem{theorem}{Theorem}}{}
\makeatother
\newcommand{\A}{\mathcal{A}}
\newcommand{\s}{\mathbf{s}}
\title[EKR-Type Theorems for Pendant Graph Constructions]{EKR-Type Theorems for Pendant Graph Constructions}
\author{Michael Carrion, Melissa M. Fuentes, Zaphenath Joseph,, Alexander Nappo}
\date{Source: arXiv:2510.22103v1 (CC BY 4.0)}
\begin{document}
\maketitle

\noindent\emph{Source and licence.} EKR-Type Theorems for Pendant Graph Constructions. Version arXiv:2510.22103v1 (CC BY 4.0), distributed under the Creative Commons Attribution 4.0 International licence, as recorded in the arXiv licence metadata for this exact version and shown on the abstract page https://arxiv.org/abs/2510.22103v1. Full rights notice: licenses/arxiv-2510.22103-CC-BY-4.0.txt.\par\smallskip

\noindent\textit{Editorial scope.} Counted unit: the Theorem whose source label is \texttt{thm:general-s} in the article (source0.tex lines 504--506, source label \texttt{thm:general-s}), delivered together with the article's own complete original proof of it (source0.tex lines 508--578, one \texttt{proof} environment of 71 lines). No mathematics is written, repaired or re-derived: statement and proof are the article's own text line for line. Required context retained uncounted: thm:KnstarEKR (lines 250--252); thm:PendantKn (lines 426--429). Every definition, display and label of those lines is retained unchanged. Omitted from this package and not redistributed: 1--249, 253--425, 430--503, 507--507, 579--769. Nothing inside a retained excerpt is omitted or altered: every retained body is delivered exactly as the article's lines are written, so no omission receipt of any kind accompanies this package. Citations: the retained excerpts carry no citation command at all, so no bibliography entry of the article is transcribed and no reference is redistributed. Reference adaptation: none was needed and none was made; every reference inside the delivered unit resolves inside it, so no unresolved reference or citation is produced. Printed slips kept literal: no printed word, symbol, hypothesis or number of the counted statement or of its proof was altered. Layout and class: the article's own class is \texttt{article} with packages including multicol, amssymb, latexsym, amsmath, amsthm, amsfonts, graphicx, tikz, enumitem, verbatim, color, xcolor, mathrsfs, ulem; the wrapper is the lane's pinned \texttt{amsart} editorial wrapper at 10pt on A4, which restates the theorem-like environments the retained text needs and adds the article's own macro definitions that the retained text needs (\texttt{\textbackslash A}, \texttt{\textbackslash s}), with extra wrapper packages (bm, bbm, dsfont, xspace, cancel, mathtools). The delivered numbering is the unit's own. Assets: no retained span contains a figure environment, a graphics-inclusion command, a PDF-page inclusion, a table or a TikZ picture; no image file is copied into this package, the package declares no asset dependency and no image file is redistributed with it. Licence: version arXiv:2510.22103v1 of this article is distributed under the Creative Commons Attribution 4.0 International licence, as recorded in the arXiv licence metadata for this exact version and shown on the abstract page https://arxiv.org/abs/2510.22103v1. A full rights notice is delivered at licenses/arxiv-2510.22103-CC-BY-4.0.txt. Characters: every retained excerpt is pure ASCII, so the delivered engine, XeLaTeX with its default Latin Modern text font, reports no missing character.
\begin{theorem}\label{thm:KnstarEKR}
For all $n \ge 2r$, the pendant complete graph $K_n^*$ is $r$-EKR. Moreover, when $n>2r$, equality holds only for an $r$-star centered at a pendant vertex.
\end{theorem}

\begin{theorem}\label{thm:PendantKn}
Let $m\ge 2$. Then the generalized pendant complete graph $K_n^m$ is $r$-EKR for all $n\ge 2r$.
In particular, the case $m=1$ is $K_n^*$ and follows from Theorem~\ref{thm:KnstarEKR}.
\end{theorem}

\begin{theorem}\label{thm:general-s}
Let $\s=(s_1,\dots,s_n)$ be a sequence of positive integers and $n\ge 2r$. Then $K_n^{\s}$ is $r$-EKR.
\end{theorem}

\begin{proof}
By relabeling the base vertices, assume $s_1\le \cdots \le s_n$, and within the largest pendant clique $K_{s_n}$ fix distinct vertices $v_1^n,v_2^n$ (if all $s_i$ are equal, then $K_n^{\s}=K_n^m$ and the result follows from Theorem~\ref{thm:PendantKn} together with the $m=1$ case from Theorem~\ref{thm:KnstarEKR}). Let $\A\subseteq \mathcal{I}^{(r)}(K_n^{\s})$ be intersecting.



Define a \emph{local} shift $T$ on sets in $\A$ by
\[
T(A)=
\begin{cases}
(A\setminus\{v_1^n\})\cup\{v_2^n\}, &\text{if } v_1^n\in A \text{ and } (A\setminus\{v_1^n\})\cup\{v_2^n\}\notin \A,\\
A,&\text{otherwise}.
\end{cases}
\]
Replacing $\A$ by $T(\A)$ if necessary, we may assume $T(\A)=\A$ (shift-stable within $K_{s_n}$).

Partition $\A$ by the presence of $v_1^n$:
\[
\A_0=\{A\in\A:\ v_1^n\notin A\},\qquad
\A_1=\{A\in\A:\ v_1^n\in A\},\qquad
\A_1'=\{A\setminus\{v_1^n\}:\ A\in\A_1\}.
\]
Let $\s'=(s_1,\dots,s_n-1)$ and $\s''=(s_1,\dots,s_{n-1})$.


\bigskip

\begin{claim}\label{claim:intersectingA1prime}
$\A_1'\subseteq \mathcal{I}^{(r-1)}\!\bigl(K_{n-1}^{\s''}\bigr)$ is intersecting.
\end{claim}
\begin{proof}
Each $A'\in \A_1'$ is an $(r-1)$-set in $K_{n-1}^{\s''}$ since we delete the single vertex $v_1^n$. 
If $A_1',A_2'\in \A_1'$ were disjoint, let $A_i=A_i'\cup\{v_1^n\}\in\A_1$ ($i=1,2$). By $T$-stability, $A_1' \cup\{v_2^n\}\in\A$. But then 
\[
\bigl(A_1'\cup\{v_2^n\}\bigr)\cap \bigl(A_2'\cup\{v_1^n\}\bigr)=\emptyset,
\]
contradicting that $\A$ is intersecting.
\end{proof}


\bigskip

\begin{claim}\label{claim:counting-star-s}
For any pendant $v_1^1$,
\[
\bigl|\mathcal{I}^{(r)}_{v_1^1}(K_n^{\s})\bigr|
\;=\;
\bigl|\mathcal{I}^{(r)}_{v_1^1}(K_n^{\s'} )\bigr|
\;+\;
\bigl|\mathcal{I}^{(r-1)}_{v_1^1}(K_{n-1}^{\s''})\bigr|.
\]
\end{claim}
\begin{proof}
Partition $\mathcal{I}^{(r)}_{v_1^1}(K_n^{\s})$ by whether $v_1^n$ is present. Removing $v_1^n$ gives a bijection from the subfamily containing $v_1^n$ to $\mathcal{I}^{(r-1)}_{v_1^1}(K_{n-1}^{\s''})$, while the subfamily avoiding $v_1^n$ is precisely $\mathcal{I}^{(r)}_{v_1^1}(K_n^{\s'})$.
\end{proof}

\bigskip

We now argue by double induction on $(r,\sum_i s_i)$ (lexicographic order). The case $r=1$ is trivial. If $s_1=\cdots=s_n$, the result follows from Theorem~\ref{thm:PendantKn} (and Theorem~\ref{thm:KnstarEKR} when $s_i=1$ for every $i$). Otherwise $s_n\ge 2$ and we apply the local shift $T$ above.

Since $\A_0$ is an intersecting family of $r$-sets in $K_n^{\s'}$ and $\A_1'$ is an intersecting family of $(r-1)$-sets in $K_{n-1}^{\s''}$ by Claim~\ref{claim:intersectingA1prime}, the inductive hypothesis gives
\[
|\A_0|\ \le\ \bigl|\mathcal{I}^{(r)}_{v_1^1}(K_n^{\s'})\bigr|\,,\qquad
|\A_1'|\ \le\ \bigl|\mathcal{I}^{(r-1)}_{v_1^1}(K_{n-1}^{\s''})\bigr|.
\]
Since $|\A|=|\A_0|+|\A_1'|$, Claim~\ref{claim:counting-star-s} yields
\[
|\A|\ \le\ \bigl|\mathcal{I}^{(r)}_{v_1^1}(K_n^{\s'})\bigr|+\bigl|\mathcal{I}^{(r-1)}_{v_1^1}(K_{n-1}^{\s''})\bigr|
\ =\ \bigl|\mathcal{I}^{(r)}_{v_1^1}(K_n^{\s})\bigr|,
\]
and hence, $K_n^{\s}$ is $r$-EKR.
\end{proof}

\end{document}
